\section{Canonical spin-two normalization}

The retained spin-two vertex has a useful entire factorization. This concerns a particular completed balanced vertex at the repeated Eisenstein parameter used in the corpus; it is not a nonvanishing theorem for every physical source functional.

\begin{theorem}[Entire amputation]\label{thm:amputation}
Let \(V(z)\) be the corpus's completed balanced spin-two vertex at repeated centered parameter \(3/2\), evaluated at projecting parameter \(iz\). Then
\begin{align}
 V(z)&=\XiC(z)^2U(z),\label{eq:amp}\\
 U(z)&=\frac{(4z^2+1)(4z^2+25)}{32}
          \XiC(z-3i)\XiC(-z-3i).\label{eq:unit}
\end{align}
The function \(U\) is entire and is the unique entire function satisfying \eqref{eq:amp}. In \(-5/2<\im z<5/2\), its zeros are precisely \(z=\pm i/2\). In particular it is nonzero on \(|\im z|<1/2\).
\end{theorem}
\cert{A2}
\begin{proof}
The completed-vertex factorization separates its projecting external square from the two shifted Xi factors and the displayed polynomial. Substitution of \(iz\) into that factorization gives \eqref{eq:amp}--\eqref{eq:unit}. Each factor in \eqref{eq:unit} is entire. In the larger open strip both shifted arguments have imaginary part less than \(-1/2\), so their Xi values are nonzero. The polynomial roots are \(\pm i/2\) and \(\pm5i/2\), with the latter two outside the open strip. For uniqueness, \(\XiC(0)\ne0\) gives a neighborhood on which cancellation of \(\XiC^2\) is valid. Any competing entire quotient agrees with \(U\) there and hence everywhere by analytic uniqueness.
\end{proof}

The external square has not disappeared from the source. Dividing by it is an analytic normalization supported by \eqref{eq:amp}; it is not permission to regard the unamputated vertex as nonzero at a Xi zero. Nor does uniqueness of this normalization construct a physical left testing field.

\section{The complete angular coefficient and forced-pole survival}

\subsection{Angular coefficient}

Denote by \(A_\rho\) the corpus's \code{physicalSelfSourceCoefficient}, obtained by assembling the residual angular period and all three principal cutoff corrections. Its normalization must be preserved: replacing that sum with only its residual part changes the following formula.

\begin{theorem}[Angular value and collision criterion]\label{thm:angular}
For a Xi zero copy with \(\re\rho>1/2\),
\begin{equation}\label{eq:fourlambda}
 A_\rho=\tfrac12\Lambda(5\rho/2-1)\Lambda(\rho/2)^2
                         \Lambda(1-3\rho/2).
\end{equation}
Moreover,
\begin{equation}\label{eq:collisions}
 A_\rho=0\quad\Longleftrightarrow\quad
 \xi(5\rho/2-1)=0\ \lor\ \xi(\rho/2)=0\ \lor\ \xi(3\rho/2)=0.
\end{equation}
\end{theorem}
\cert{A3}\cert{A4}
\begin{proof}
The three principal Mellin terms sum to the native principal moment. Each cutoff quotient has exactly the scalar factor needed to cancel the principal correction in the residual period. Substituting the evaluated residual period gives \eqref{eq:fourlambda}; this is the content of the source's cutoff-assembly identity, not an appeal to formal integration of a divergent series. All three affine arguments are nonreal because \(\im\rho\ne0\). Their Gamma and elementary pole-clearing factors are finite and nonzero, so vanishing of \(\Lambda\) is equivalent to vanishing of \(\xi\) at each argument. Applying \(\xi(1-u)=\xi(u)\) to the last factor gives \eqref{eq:collisions}.
\end{proof}

Equation \eqref{eq:fourlambda} evaluates a source; it does not prove that source nonzero. Its strict-right restriction is also essential to its present formal application. The value of this angular functional and the entire source \(U\) in Theorem~\ref{thm:amputation} are different objects.

\subsection{Onset and the finite-part radial moment}

Let \(\chi\) be the fixed smooth transition used in the corpus, zero before 0 and one after 1. Its derivative has integral one. Define
\[
 O(s,w)=\int_0^1\chi'(u)e^{(2s-w)u}\dd u,\qquad
 O_L(s,w)=\int_0^1\chi'(u)e^{L(2s-w)u}\dd u.
\]
These are \code{clearedOnset} and \code{clearedOnsetDuration}; the second formula stretches the exponent rather than silently replacing an already constructed solution by another one.

For \(0\le z<1\), define
\begin{align*}
 c_0(s)&=(5s^2+15s+4)/4,\\
 c_1(s,w)&=\frac{s(s-1)}4\left((s+2)^2+(s+1)(s+2)
                                  -\frac{(w+1)(w+2)}2\right),\\
 F(w,s,a;z)&=a^2\sqrt{1-z}\,(1+\sqrt{1-z})^{2s-w}
                      [c_0(s)+c_1(s,w)(1-z)].
\end{align*}
Powers of a positive real base use its real logarithm. Write \(b=1/2\). The radial moment used here is precisely
\begin{align}\label{eq:moment}
 M(w,s,a)={}&\frac{F(w,s,a;0)b^{w+1}}{w+1}
 +\int_0^b z^w[F(w,s,a;z)-F(w,s,a;0)]\dd z\nonumber\\
 &+\int_b^1z^wF(w,s,a;z)\dd z.
\end{align}
This is the source definition \code{fundamentalSourceMoment}. For the physical nonreal frequencies in \eqref{eq:testing}, \(\re w>-2\) and \(w\ne-1\); the endpoint subtraction gains the one power needed at zero. Formula \eqref{eq:moment} must not be replaced by an unsubtracted improper integral in that chamber.

\begin{theorem}[A common nonvanishing onset duration]\label{thm:onset}
There is one \(L>0\) such that \(O_L(s_b(\rho),w_e(\rho))\ne0\) for every Xi zero copy, both parent branches \(b\), and both testing sides \(e\). Also \(O(s,2s)=1\) and \(O_1(s,w)=O(s,w)\).
\end{theorem}
\cert{A6}\cert{A7}
\begin{proof}
The pulse transform \(z\mapsto\int_0^1\chi'(u)e^{zu}\dd u\) is entire and takes value one at zero. For each fixed pair \((s,w)\), its composition with \(L\mapsto L(2s-w)\) is entire and not identically zero. Its real nonzero set is open and dense; if \(2s=w\), it is the whole real line. Xi zero copies form a countable type, and the two branch choices are finite. Baire's theorem makes the intersection of these open dense sets dense. It meets \((0,\infty)\). The two identities follow by substitution and the normalization of \(\chi'\).
\end{proof}

This theorem establishes no preferred numerical duration and, in particular, does not certify \(L=1\) at every target frequency. Transporting the complete source, response, boundary terms, and time convention to the chosen \(L\) remains an additional obligation.

\subsection{Exact residue factorization}

Let \(R_{q,b}(w;a,x)\) denote \code{continuedPhysicalSelfResponse}, with tortoise observation point \(x<0\). Set \(z(x)=\tanh^2x\), and
\[
 m_w(x)=z(x)^{w/2}(1-z(x))^{1/2}.
\]
Let \(D'_{q,b,e}\) be \code{physicalFundamentalDenominatorSlope}. The corpus's simple-connection-zero calculation makes it nonzero. It includes both the connection slope and the source-rate clearing factor; replacing it by an unnamed Wronskian constant would lose its normalization.

\begin{theorem}[Five-factor survival criterion]\label{thm:survival}
At \(w=w_e(\rho)\), the physical Laurent coefficient is
\begin{equation}\label{eq:residue}
 \Res R_{q,b}(\,\cdot\,;a,x)
 =m_w(x)\frac{A_\rho M(w,s_b,a)O(s_b,w)}{D'_{q,b,e}}.
\end{equation}
For \(\re\rho>1/2\), \(R_{q,b}\) has no finite punctured-neighborhood limit at \(w_e\) if and only if
\begin{equation}\label{eq:five}
 \begin{gathered}
 \xi(5\rho/2-1)\ne0,\qquad\xi(\rho/2)\ne0,\qquad\xi(3\rho/2)\ne0,\\
 M(w_e,s_b,a)\ne0,\qquad O(s_b,w_e)\ne0.
 \end{gathered}
\end{equation}
\end{theorem}
\cert{A5}
\begin{proof}
At a fundamental connection zero, the two endpoint source integrals combine into the elementary nonzero mode \(m_w\) times a single moment. To see why the moment is independent of the interior splitting point, differentiate their sum: the derivative of the lower finite-part integral is the original integrand, and the derivative of the upper integral is its negative. The sum is constant and equals \eqref{eq:moment} at \(b=1/2\). The cleared Green numerator therefore factors as \(m_wMO\). Dividing by the simple denominator zero gives \eqref{eq:residue}. On \(x<0\), both bases in \(m_w\) are positive and nonzero; the denominator slope is nonzero. A simple-pole quotient has a nonremovable singularity exactly when its numerator at the pole is nonzero. Theorem~\ref{thm:angular} then yields \eqref{eq:five}.
\end{proof}

The word ``criterion'' is deliberate. No universal nonvanishing of the last two factors is asserted. The theorem allows \(a=0\), in which case the moment vanishes; a discussion of excitation must first choose a nonzero amplitude. A nonzero radial norm, evaluated next, is not a nonzero source moment.

\section{Radial regularization and Laurent extraction}

The mode-square finite part in the corpus is
\begin{equation}\label{eq:betanorm}
 N(w)=\frac{\Gamma(w)\Gamma(3/2)}{2\Gamma(w+3/2)},
 \qquad \re w>-2,\quad\im w\ne0.
\end{equation}
It is nonzero on that domain. This is the continued, horizon-subtracted mode square, not a positive Hilbert-space norm.
\cert{A8}
Its proof applies the endpoint-subtracted beta identity with parameters \(w,3/2\). The nonreal restrictions exclude all Gamma poles, and Gamma has no zeros, so numerator and denominator are nonzero. The analytic continuation and endpoint subtraction are part of the formal norm theorem.

There is also an actual integral--Laurent interchange for the horizon-mode temporal current. With \(w_0=w_e(\rho)\), let \(J(u,x)\) denote that current multiplied by its radius Jacobian, and let \(D(u)\) be its clearing factor. The theorem constructs analytic functions \(c_0(u),c_1(u)\) and a meromorphic function \(\Phi(u)\) such that, eventually for \(u\ne w_0\),
\begin{equation}\label{eq:radiallimit}
 \lim_{X\to-\infty}\left\{\int_X^0J(u,x)\dd x+
 \frac{c_0(u)z(X)^u/u+c_1(u)z(X)^{u+1}/(u+1)}{D(u)}\right\}
 =\Phi(u).
\end{equation}
Its Laurent coefficient of order \(-1\) is exactly the independently defined \code{physicalRadialResidueFinitePart}.
\cert{A9}

The proof first establishes a frequency-uniform second-order remainder bound for the cleared horizon density. Subtraction of the two displayed jets leaves an integrable Mellin remainder on the target frequency neighborhood. The complementary compact integral is analytic by domination. Factoring \(D(u)=(u-w_0)\kappa(u)\) with \(\kappa(w_0)\ne0\) identifies the residue from the cleared density at \(w_0\); uniqueness of its jet coefficients matches the prescribed counterterms. This argument justifies the interchange for this current. It gives meromorphy with a possible simple pole; it does not assert a nonzero residue when the source factor vanishes.

The branch selector in this result is \code{outer = false}. The corresponding normalizable-mode leg, all angular integration, and limits in time or the null horizon are not conclusions of \eqref{eq:radiallimit}. Iterating these operations requires its own uniform estimates and identification of common counterterms.
